\documentclass[12pt, a4paper]{article}
\usepackage[utf8]{inputenc}
\usepackage[T1]{fontenc}
\usepackage[french]{babel}
\usepackage[margin=2.5cm]{geometry}
\usepackage{amsmath, amssymb, amsthm, mathtools}
\usepackage{listings}
\usepackage{xcolor}
\usepackage{tcolorbox}
\tcbuselibrary{theorems, skins, breakable}
\usepackage{booktabs}
\usepackage{fancyhdr}
\usepackage{hyperref}
\usepackage{tikz}
\usetikzlibrary{arrows.meta, positioning, shapes.geometric}
\usepackage{algorithm}
\usepackage{algpseudocode}

\definecolor{suva_blue}{RGB}{30, 100, 180}
\definecolor{code_bg}{RGB}{245, 247, 250}
\definecolor{code_kw}{RGB}{0, 80, 200}
\definecolor{code_str}{RGB}{180, 40, 40}
\definecolor{code_cmt}{RGB}{80, 130, 80}
\definecolor{warn_bg}{RGB}{255, 248, 220}
\definecolor{success_bg}{RGB}{230, 255, 235}

\lstdefinestyle{python}{
  language=Python,
  basicstyle=\ttfamily\small,
  keywordstyle=\color{code_kw}\bfseries,
  stringstyle=\color{code_str},
  commentstyle=\color{code_cmt}\itshape,
  backgroundcolor=\color{code_bg},
  frame=single, rulecolor=\color{gray!40},
  numbers=left, numberstyle=\tiny\color{gray},
  breaklines=true, tabsize=4, showstringspaces=false,
}
\lstdefinestyle{output}{
  basicstyle=\ttfamily\small,
  backgroundcolor=\color{gray!10},
  frame=single, rulecolor=\color{gray!50},
  breaklines=true,
}

\newtcolorbox{defbox}[2][]{colback=success_bg,colframe=green!50!black,
  fonttitle=\bfseries,title={Définition : #2},#1}
\newtcolorbox{theorembox}[2][]{colback=blue!5,colframe=suva_blue,
  fonttitle=\bfseries,title={Théorème : #2},#1}
\newtcolorbox{exercice}[2][]{colback=white,colframe=suva_blue!60,
  fonttitle=\bfseries,title={Exercice #2},#1,breakable}
\newtcolorbox{solution}[1][]{colback=gray!5,colframe=gray!40,
  fonttitle=\bfseries\itshape,title={Solution},#1,breakable}
\newtcolorbox{importantbox}[1][]{colback=suva_blue!10,colframe=suva_blue,
  fonttitle=\bfseries,title={Connexion avec le projet},#1}
\newtcolorbox{warningbox}[1][]{colback=warn_bg,colframe=orange!80!black,
  fonttitle=\bfseries,title={Attention},#1}

\newcommand{\Zq}{\mathbb{Z}_q}
\newcommand{\Rq}{R_q}

\pagestyle{fancy}
\fancyhf{}
\fancyhead[L]{\color{suva_blue}\textbf{SuVa --- Chapitre 5}}
\fancyhead[R]{\color{gray}Dilithium Complet}
\fancyfoot[L]{\footnotesize\color{gray}Mohamed Lamine MBENGUE}
\fancyfoot[C]{\thepage}
\fancyfoot[R]{\footnotesize\color{gray}portfolio-alamine.web.app}
\renewcommand{\headrulewidth}{0.4pt}
\renewcommand{\footrulewidth}{0.3pt}

% ══════════════════════════════════════════════════════════════════════
\begin{document}

\begin{center}
  {\color{suva_blue}\rule{\linewidth}{2pt}}\\[0.5cm]
  {\Huge\bfseries\color{suva_blue} Chapitre 5}\\[0.3cm]
  {\large\bfseries Dilithium / ML-DSA --- Le schéma de signature complet}\\[0.2cm]
  {\small\color{gray} Keygen, Sign, Verify : code source et mathématiques détaillés}\\[0.3cm]
  {\color{suva_blue}\rule{\linewidth}{2pt}}
\end{center}

\vspace{0.3cm}
\begin{tcolorbox}[colback=suva_blue!5, colframe=suva_blue!40,
  left=6pt, right=6pt, top=4pt, bottom=4pt]
\small
\begin{tabular}{ll}
  \textbf{Auteur}      & Mohamed Lamine MBENGUE \\
  \textbf{Spécialités} & Sécurité · DevOps · Dev Full Stack Web \& Mobile \\
  \textbf{Contact}     & +221 78 178 44 36 · mbenguemohamedlamine2022@gmail.com \\
  \textbf{Portfolio}   & \href{https://portfolio-alamine.web.app}{\texttt{portfolio-alamine.web.app}} \\
\end{tabular}
\end{tcolorbox}
\vspace{0.3cm}

\tableofcontents
\newpage

%══════════════════════════════════════════════════════════════════════
\section{Vue d'ensemble du schéma Dilithium}
%══════════════════════════════════════════════════════════════════════

\subsection{Paramètres Dilithium2 (niveau utilisé dans SuVa)}

\begin{center}
\begin{tabular}{clp{7cm}}
\toprule
\textbf{Param.} & \textbf{Valeur} & \textbf{Rôle} \\
\midrule
$n$    & 256 & Degré des polynômes \\
$q$    & 8\,380\,417 & Module premier \\
$k$    & 4 & Lignes de la matrice $A$ \\
$l$    & 4 & Colonnes de la matrice $A$ \\
$\eta$ & 2 & Borne des coefficients secrets ($s_1, s_2$) \\
$\tau$ & 39 & Nb de $\pm 1$ dans le challenge $c$ \\
$\gamma_1$ & $2^{17} = 131\,072$ & Plage de masquage de $z$ \\
$\gamma_2$ & $(q-1)/88 = 95\,232$ & Plage de rounding \\
$d$    & 13 & Bits droppés de la clé publique \\
$\omega$ & 80 & Max de 1s dans le hint $h$ \\
$\beta$ & 78 & Borne de rejet ($= \tau \times \eta$) \\
\bottomrule
\end{tabular}
\end{center}

\subsection{Les trois algorithmes}

\begin{center}
\begin{tikzpicture}[
  block/.style={draw=suva_blue, rounded corners, fill=blue!5,
    minimum width=5cm, minimum height=1.2cm,
    text width=4.8cm, align=center},
  arrow/.style={->, thick, suva_blue},
  label/.style={font=\small, color=gray}
]
  \node[block] (kg) {
    \textbf{KeyGen} (off-chain, une fois)\\
    \footnotesize $\zeta \to (A, s_1, s_2) \to t = As_1+s_2$\\
    \footnotesize \texttt{pk = (rho, t1), sk = (rho, K, tr, s1, s2, t0)}
  };
  \node[block, below=1.5cm of kg] (sg) {
    \textbf{Sign} (off-chain, par message)\\
    \footnotesize Boucle de rejet jusqu'à trouver $(z, h)$ valide\\
    \footnotesize \texttt{sig = (c\_tilde, z, h)}
  };
  \node[block, below=1.5cm of sg] (vf) {
    \textbf{Verify} (on-chain Solidity)\\
    \footnotesize Recalcule $w' = Az - ct_1 \cdot 2^d$\\
    \footnotesize Compare $\tilde{c} \stackrel{?}{=} H(\mu \| w')$
  };
  \draw[arrow] (kg) -- node[label, right] {sk → Sign} (sg);
  \draw[arrow] (sg) -- node[label, right] {(pk, sig, msg)} (vf);
  \node[right=2cm of kg, font=\scriptsize, align=left, color=gray] {
    Fichier :\\
    \texttt{Falcon\_dilithium.py}
  };
  \node[right=2cm of vf, font=\scriptsize, align=left, color=gray] {
    Fichier :\\
    \texttt{ZKNOX\_dilithium.sol}
  };
\end{tikzpicture}
\end{center}

%══════════════════════════════════════════════════════════════════════
\section{KeyGen --- Génération de clés en détail}
%══════════════════════════════════════════════════════════════════════

\subsection{Algorithme}

\begin{lstlisting}[style=python]
# Extrait de Falcon_dilithium/Falcon_dilithium.py - methode keygen()

def keygen(self, _xof=None, _xof2=None):
    # 1. Graine principale : 32 bytes aleatoires
    zeta = self.random_bytes(32)

    # 2. Expander : SHAKE256(zeta) -> 128 bytes
    #    Divise en rho (32), rho_prime (64), K (32)
    seed_bytes = self._h(zeta, 128)
    rho        = seed_bytes[:32]   # seed pour matrice publique A
    rho_prime  = seed_bytes[32:96] # seed pour vecteurs secrets s1, s2
    K          = seed_bytes[96:]   # seed pour la signature (nonce)

    # 3. Generer la matrice publique A in R_q^{k x l}
    # A est generee par rejection sampling depuis SHAKE128(rho)
    # IMPORTANTE : A est dans le domaine NTT deja (A_hat)
    A_hat = self._expand_matrix_from_seed(rho)
    # A_hat est k x l = 4x4 matrices de polynomes

    # 4. Generer les vecteurs secrets
    # s1 in R_q^l : coefficients dans [-eta, eta]
    # s2 in R_q^k : coefficients dans [-eta, eta]
    s1, s2 = self._expand_vector_from_seed(rho_prime)
    # Pour Dilithium2 (eta=2) : coefficients dans {-2,-1,0,1,2}

    # 5. Calculer t = A * s1 + s2
    t = (A_hat @ s1.to_ntt()).from_ntt() + s2
    # Decomposition : t = t1 * 2^d + t0
    t1, t0 = t.power_2_round(self.d)

    # 6. Construire les cles
    pk = self._pack_pk(rho, t1)    # 1312 bytes
    tr = self._h(pk, 32)           # Hash de la cle publique
    sk = self._pack_sk(rho, K, tr, s1, s2, t0)  # 2528 bytes

    return pk, sk
\end{lstlisting}

\subsection{Tailles}

\begin{center}
\begin{tabular}{lll}
\toprule
\textbf{Composant} & \textbf{Taille} & \textbf{Contenu} \\
\midrule
$\rho$ & 32 bytes & Graine pour générer $A$ \\
$t_1$ & 1280 bytes & $k \times n \times 10$ bits $= 4 \times 256 \times 10 / 8$ \\
\textbf{Total pk} & \textbf{1312 bytes} & \\
\midrule
$\rho$ & 32 bytes & \\
$K$ & 32 bytes & \\
$\text{tr}$ & 32 bytes & Hash de pk \\
$s_1$ & 384 bytes & $l \times n \times (2\eta+1)$-bits codage \\
$s_2$ & 384 bytes & $k \times n \times (2\eta+1)$-bits codage \\
$t_0$ & 1664 bytes & $k \times n \times 13$ bits \\
\textbf{Total sk} & \textbf{2528 bytes} & \\
\bottomrule
\end{tabular}
\end{center}

%══════════════════════════════════════════════════════════════════════
\section{Sign --- La boucle de rejet}
%══════════════════════════════════════════════════════════════════════

\begin{warningbox}
La boucle de rejet est le cœur de Dilithium.
En moyenne, \textbf{4.25 itérations} sont nécessaires pour Dilithium2.
\end{warningbox}

\begin{lstlisting}[style=python]
# Extrait simplifie de sign()

def sign(self, sk_bytes, m, _xof=None, _xof2=None):
    rho, K, tr, s1, s2, t0 = self._unpack_sk(sk_bytes)
    A_hat = self._expand_matrix_from_seed(rho)

    # Message representative : mu = SHAKE256(tr || m, 64 bytes)
    mu = self._h(tr + m, 64)

    # rho_prime pour le masque y
    rho_prime = self._h(K + mu, 64)

    # Precomputer en NTT (important pour la performance)
    s1 = s1.to_ntt()
    s2 = s2.to_ntt()
    t0 = t0.to_ntt()

    alpha = 2 * self.gamma_2
    kappa = 0  # nonce incremente a chaque iteration

    while True:  # Boucle de rejet
        # Etape 1 : Tirer le masque y
        # y in R_q^l avec ||y||_inf <= gamma_1
        y = self._expand_mask_vector(rho_prime, kappa)
        kappa += self.l

        # Etape 2 : Calculer w = A*y dans R_q^k
        w = (A_hat @ y.to_ntt()).from_ntt()

        # Etape 3 : Decomposer w en w1 (haute) et w0 (basse)
        w1, w0 = w.decompose(alpha)

        # Etape 4 : Hacher pour obtenir le challenge
        # c_tilde = H(mu || w1_packed)
        c_tilde = self._h(mu + w1.bit_pack_w(self.gamma_2), 32)
        c = self.R.sample_in_ball(c_tilde, self.tau).to_ntt()

        # Etape 5 : z = y + c*s1
        z = y + (s1.scale(c)).from_ntt()

        # REJET 1 : z trop grand -> fuit info sur s1
        if z.check_norm_bound(self.gamma_1 - self.beta):
            continue  # Recommencer

        # Etape 6 : Calculer les termes de correction
        w0_minus_cs2 = w0 - s2.scale(c).from_ntt()

        # REJET 2 : w0 - cs2 trop grand
        if w0_minus_cs2.check_norm_bound(self.gamma_2 - self.beta):
            continue

        c_t0 = t0.scale(c).from_ntt()

        # REJET 3 : c*t0 trop grand
        if c_t0.check_norm_bound(self.gamma_2):
            continue

        # Etape 7 : Construire le hint
        h = (w0_minus_cs2 + c_t0).make_hint_optimised(w1, alpha)

        # REJET 4 : trop de 1s dans h (trop grande correction)
        if h.sum_hint() > self.omega:
            continue

        # Succes ! Empaqueter la signature
        return self._pack_sig(c_tilde, z, h)
        # Signature = (c_tilde, z, h) = (32, ~1152, ~100) bytes
\end{lstlisting}

\subsection{Pourquoi le rejet protège la clé secrète ?}

\begin{theorembox}{Rejet pour la confidentialité}
Si on publiait $z = y + cs_1$ sans restriction, un attaquant qui observe
plusieurs signatures pourrait faire de l'analyse statistique~:

$z = y + cs_1 \Rightarrow z - y = cs_1$

Si la distribution de $z$ dépend de $s_1$, la clé secrète fuit.

\textbf{Le rejet uniformise~:} En rejetant les $z$ trop grands (proches
de la borne), on s'assure que la distribution de $z$ est essentiellement
\textbf{uniforme} sur $[-\gamma_1+\beta, \gamma_1-\beta]$ et n'informe
pas sur $s_1$.
\end{theorembox}

%══════════════════════════════════════════════════════════════════════
\section{Verify --- Vérification complète}
%══════════════════════════════════════════════════════════════════════

\begin{lstlisting}[style=python]
# Extrait simplifie de verify()

def verify(self, pk_bytes, m, sig_bytes, _xof=None):
    # Decompresser
    rho, t1 = self._unpack_pk(pk_bytes)
    c_tilde, z, h = self._unpack_sig(sig_bytes)

    # Verifications basiques
    if h.sum_hint() > self.omega:
        return False
    if z.check_norm_bound(self.gamma_1 - self.beta):
        return False

    # Regenerer A (depuis la graine rho)
    A_hat = self._expand_matrix_from_seed(rho)

    # Recalculer mu
    tr = self._h(pk_bytes, 32)
    mu = self._h(tr + m, 64)

    # Reconstruire c depuis c_tilde
    c = self.R.sample_in_ball(c_tilde, self.tau).to_ntt()

    # Calculer Az - c*t1*2^d
    # t1 * 2^d reconstruit le vecteur t1 * 2^d dans Z_q
    z_ntt = z.to_ntt()
    t1_scaled = t1.scale(1 << self.d).to_ntt()
    Az_minus_ct1 = ((A_hat @ z_ntt) - t1_scaled.scale(c)).from_ntt()
    # Note : (A*z - c*t1*2^d) ≈ A*y (si z = y + c*s1 et t1*2^d ≈ t-t0)

    # Utiliser le hint pour recuperer w1'
    w_prime = h.use_hint(Az_minus_ct1, 2 * self.gamma_2)

    # Verification finale
    c_tilde_check = self._h(mu + w_prime.bit_pack_w(self.gamma_2), 32)
    return c_tilde == c_tilde_check
\end{lstlisting}

\subsection{Preuve de correction (intuition)}

Lors de la signature on a~:
$$z = y + cs_1, \quad t = As_1 + s_2, \quad t = t_1 \cdot 2^d + t_0$$

Lors de la vérification~:
\begin{align*}
Az - ct_1 \cdot 2^d
&= A(y + cs_1) - ct_1 \cdot 2^d \\
&= Ay + cAs_1 - ct_1 \cdot 2^d \\
&= Ay + c(t - s_2) - c(t - t_0) \\
&= Ay - cs_2 + ct_0 \\
&\approx Ay \quad \text{(si } \|cs_2\|_\infty, \|ct_0\|_\infty < \gamma_2\text{)}
\end{align*}

Le hint corrige la petite différence entre $Ay$ et $Az - ct_1 \cdot 2^d$
pour retrouver $w_1 = \text{HighBits}(Ay)$.

%══════════════════════════════════════════════════════════════════════
\section{ETHDilithium --- les deux changements clés}
%══════════════════════════════════════════════════════════════════════

\subsection{Changement 1 : SHAKE → Keccak256}

\begin{center}
\begin{tabular}{lll}
\toprule
& \textbf{Dilithium standard} & \textbf{ETHDilithium} \\
\midrule
Hash function & SHAKE-128/256 & Keccak-256 \\
Implémentation Solidity & 269 lignes & opcode natif (3 gaz) \\
Coût on-chain & $\sim 50\%$ du gas total & $\sim 5\%$ \\
Conformité NIST & Oui (FIPS 204) & Non \\
\bottomrule
\end{tabular}
\end{center}

\subsection{Changement 2 : Clé publique étendue}

\begin{center}
\begin{tabular}{lll}
\toprule
& \textbf{Dilithium} & \textbf{ETHDilithium} \\
\midrule
Clé publique & $(\rho, t_1)$ & $(\hat{A}, t_1)$ \\
Taille pk & 1312 bytes & $\sim 4096 + 1280$ bytes \\
Coût keygen & Normal & On précalcule $\hat{A}$ \\
Coût verify & Doit régénérer $\hat{A}$ & $\hat{A}$ déjà là \\
Économie gas & --- & $\sim -50\%$ \\
\bottomrule
\end{tabular}
\end{center}

\begin{lstlisting}[style=python]
# Extrait de Falcon_dilithium_eth.py

def keygen(self, _xof=None, _xof2=None):
    # ... (pareil jusqu'a t1, t0)

    # DIFFERENCE : stocker A_hat dans la cle publique
    # Au lieu de (rho, t1), on stocke (A_hat, t1_new)
    pk_new = self._pack_pk_new(A_hat, t1)
    # pk_new est plus grande (toute la matrice A)
    # mais la verification n'a plus besoin de regenerer A

    # t1_new : deja en NTT form dans la pk
    # Pour que la verification puisse directement utiliser
    # t1 * 2^d sans recalculer le NTT
    return pk_new, sk
\end{lstlisting}

%══════════════════════════════════════════════════════════════════════
\section{Exercices}
%══════════════════════════════════════════════════════════════════════

\begin{exercice}{1 --- Tracer le flow complet}
Trace sur papier le flow complet de Dilithium en annotant chaque étape
avec le fichier/fonction correspondant dans le projet~:

\begin{enumerate}[label=\alph*)]
  \item KeyGen~: de \texttt{random\_bytes(32)} jusqu'à \texttt{pack\_pk}
  \item Sign~: de \texttt{\_unpack\_sk} jusqu'à \texttt{return \_pack\_sig}
  \item Verify~: de \texttt{\_unpack\_pk} jusqu'au \texttt{return True/False}
\end{enumerate}
\end{exercice}

\begin{exercice}{2 --- Pourquoi le nonce kappa}
Dans \texttt{sign()}, \texttt{kappa} est incrémenté de $l$ à chaque itération.
\begin{enumerate}[label=\alph*)]
  \item Pourquoi ne pas utiliser \texttt{os.urandom()} pour le masque $y$ ?
  \item Que se passerait-il si deux signatures utilisaient le même $y$ ?
  \item Comment \texttt{kappa} garantit-il que $y$ est différent à chaque itération ?
\end{enumerate}
\end{exercice}

\begin{solution}
\begin{lstlisting}[style=python]
# b) Si y est reutilise pour deux messages m1, m2 :
# sig1 = (c1, z1=y+c1*s1, h1)
# sig2 = (c2, z2=y+c2*s1, h2)
# Alors : z1 - z2 = (c1 - c2) * s1
# Si c1 != c2, on peut isoler s1 ! -> CATASTROPHIQUE

# c) kappa est incremente : rho_prime || kappa
# -> deux iterations differentes -> y differents
# Le PRNG est deterministe depuis (rho_prime, kappa)
# -> securite ET determinisme (utile pour les KAT)
\end{lstlisting}
\end{solution}

\begin{exercice}{3 --- Sécurité du hint}
Le hint $h$ peut contenir au plus $\omega = 80$ valeurs 1.
\begin{enumerate}[label=\alph*)]
  \item Combien de bits d'information le hint encode-t-il au maximum ?
        ($k \times n$ bits potentiels, mais $\leq \omega$ sont 1)
  \item Pourquoi la limite $\omega$ est-elle nécessaire pour la sécurité ?
  \item Que se passe-t-il si un attaquant peut forcer $\|h\|_1 > \omega$ ?
\end{enumerate}
\end{exercice}

\begin{exercice}{4 --- ETHDilithium vs Dilithium en Python}
\begin{enumerate}[label=\alph*)]
  \item Ouvre \texttt{example.py} et \texttt{example\_zk.py}.
  \item Lance les deux exemples et observe la sortie.
  \item Mesure et compare les temps d'exécution de ETHDilithium2 vs Dilithium2
        pour keygen, sign et verify.
  \item Pourquoi ETHDilithium est plus lent en Python mais plus rapide on-chain ?
\end{enumerate}
\end{exercice}

\end{document}
